\documentclass[11pt,a4paper]{article}
\usepackage[T1]{fontenc}
\usepackage{lmodern,amsmath,amssymb,amsthm,mathtools}
\usepackage[margin=26mm]{geometry}
\usepackage{array,longtable,booktabs,xcolor,microtype}
\usepackage[colorlinks=true,linkcolor=blue!45!black,citecolor=blue!45!black,urlcolor=blue!45!black]{hyperref}
\newtheorem{theorem}{Theorem}[section]
\newtheorem{lemma}[theorem]{Lemma}
\newtheorem{proposition}[theorem]{Proposition}
\newtheorem{corollary}[theorem]{Corollary}
\theoremstyle{definition}
\newtheorem{definition}[theorem]{Definition}
\theoremstyle{remark}
\newtheorem{remark}[theorem]{Remark}
\newcommand{\Q}{\mathbb Q}
\newcommand{\Z}{\mathbb Z}
\newcommand{\Sha}{\operatorname{Sha}}
\setlength{\emergencystretch}{2em}
\setcounter{tocdepth}{1}
\title{An explicit quartic construction over the rationals}
\author{}
\date{7 November 2026}
\begin{document}
\maketitle
\begin{abstract}
We give a complete quartic coefficient table in 456 total
variables, obtained from an explicit quadratic arithmetic system.
The algebraic transformations, the local normalization, the index-depth
calculation and the coefficient-height formula are proved. The rational
input projection is the integers under four stated source arithmetic
inputs and an explicitly specified weighted prime-selection input.
The polynomial has primitive integer coefficients and height
$64\cdot3^{105665812}$. Complete data and independent reconstruction
checks accompany the paper.
\end{abstract}
\tableofcontents
\newpage
\section{An explicit quartic and its meaning}
A polynomial $Q\in\mathbb Z[T,Y_1,\ldots,Y_m]$ has rational
input projection $\{t\in\mathbb Q:\exists\mathbf y\in\mathbb Q^m,
Q(t,\mathbf y)=0\}$. All variable counts include the input $T$.
The coefficient height is the largest absolute coefficient.

The supplied polynomial is the sum of the squares of a fully listed
quadratic system. Its coefficient table, degree, primitivity and
height are exact algebraic assertions. The equality of its input
projection with $\mathbb Z$ uses the arithmetic inputs of
Section~\ref{sec:u29} and the weighted prime-selection input of
Section~\ref{sec:asymmetric-depth}. The constructions below prove
the transformations and their soundness and completeness on the
stated domains; the source arithmetic premises remain explicit.

The starting arithmetic system, its local extraction argument and
the simultaneous approximation construction are given in
\cite{SmallVariables2026,RepositoryV9}. The present paper develops
the quadratic implementation, its exact local index conditions,
and the transformations needed for the supplied coefficient table.
\section{Arithmetic inputs and the product-height identity}\label{sec:u29}
\subsection{Arithmetic inputs and conventions}

Put
\[
 K=\mathbb Q(\sqrt2),\quad d=3-\sqrt2,\quad D=d^2=11-6\sqrt2,
 \qquad E:y^2=x^3-Dx,
\]
and fix
\[
 R_0=(1+\sqrt2,2\sqrt2),\qquad
 P=2R_0=\left(\frac{57}{8}-\frac72\sqrt2,
          -\frac{165}{8}+\frac{417}{32}\sqrt2\right).
\]
Write $\sigma$ for conjugation, $T_0=(0,0)$, $f(Q)=1/x(Q)$,
and $z(Q)=-x(Q)/y(Q)$.  Heights are absolute logarithmic Weil
heights, with normalized local degree weights.  Valuations at
unramified primes satisfy $v(p)=1$.

We use the following arithmetic hypotheses, with the named source
locations in \cite{RepositoryV9}.
\begin{enumerate}
\item (Elliptic group and height.)\\
      Source: \nolinkurl{subsec:v5-point-source}, \nolinkurl{eq:o8-indexheight}.\par $2E(K)=\mathbb ZP$, and
      $h(x([n]P))\ge n^2-6$ for every nonzero integer $n$.
\item (Contact height: Theorem \nolinkurl{thm:quartic-height19}.)
      For $t=u/v$ in lowest terms, $v>0$, and
      $t\notin\mathcal D=\{0,1,2401,-343/32\}$, define $M_5(t)$
      to be the product of the primes $p>76176$ for which at least one of
\[
 v_p(v),\quad v_p(u),\quad v_p(u-v),\quad
 v_p(u-2401v),\quad v_p(32u+343v)
\]
      is positive and odd.  These are the normalized contacts with
      $\infty,0,1,2401,-343/32$.  Put
\[
 B_* =\prod_{p\le76176}p,\qquad
 H=H_{19}=2^{110}210^{1800}B_*^{19}+128.
\]
      The inherited assertion is $h(t)\le H M_5(t)^{19}$.
      Exact integer evaluation gives $H<2^{2093448}$.

\item (Pole obstruction and realization.)\\
      Source: \nolinkurl{hyp:o8-pole} and
      \nolinkurl{lem:original-seven-simultaneous-realization}.
      The pole equations, with $B=17b$,
\[
 U=b\frac{1-bt^2}{1+bt^2},\quad h_0=e(B-U^2),\quad
 p^2=B(1-e)(1+3e),\quad
 Bq^2=(s^2-Bh_0)(s^2+3Bh_0),
\]
      exclude negative odd $v_\ell(b)$ for $\ell\notin\{2,3,17\}$.
      They are soluble with $t=0$ whenever $b$ has square reduced
      denominator, is positive and less than $17$, is square in
      $\mathbb Q_2,\mathbb Q_3,\mathbb Q_{17}$, and satisfies
      $v_2(b)\ge3$, $v_3(b)\ge1$, $v_{17}(b)\ge2$.
\item (Simultaneous prime patterns: Theorem \nolinkurl{thm:v5-simultaneous-primes},
      Lemma \nolinkurl{lem:v5-two-reductions}, and the application of
      \cite[Theorem 4.1]{Sun2026Rational86}.)
      The simultaneous Gaussian prime-pattern input supplies the
      prescribed prime values $u+c_jw^2$, arbitrary fixed divisibility
      of $w$, and arbitrarily large $u/w^2$.  Their two conjugate
      reduction orders satisfy
\[
 v_2(d_-)=6,\quad v_2(d_+)=2,\quad
 \gcd(d_-(p),d_+(p'))=4
\]
      for every pair of selected primes.  An arbitrary fixed odd
      avoidance modulus can be imposed.  For each selected Gaussian
      prime $\pi$, the source convention is
      $\pi\equiv1+3i\pmod7$, making both $N(\pi-1)$ and
      $N(\pi+1)$ units at seven.  At each other odd prime $\ell$
      in the fixed avoidance modulus, the condition
      $\pi\equiv i\pmod\ell$ makes both orders congruent to two
      modulo $\ell$.  Thus the orders avoid the entire odd part
      of the fixed modulus.  The selected rational primes are $1$
      modulo $8\cdot3\cdot17$.
\end{enumerate}
These hypotheses include the simultaneous application of Sun's stated
prime-pattern theorem.  They are assumptions of the theorem below,
not conclusions of its implementation audit.
The local depth of the displayed point is proved in Proposition~\ref{anchor3:exact-depth}.

\subsection{An elementary Lucas growth estimate}

The following estimate is independent of the arithmetic hypotheses.
It supplies the growth inequality used in the integer-extraction argument.

\begin{lemma}\label{lem:lucas-growth-elementary}
Let $a,t,b\in\mathbb Z$, with $|a|>7$ and $b\ne0$, satisfy
\[
 t^2-a(a-2)tb+(a-2)^2b^2=1.
\]
Then $\log|t|>|a|-3$. In particular,
$\log\max(|a|,|t|,|b|)>|a|-3$.
\end{lemma}
\begin{proof}
Put $A=|a|$, $c=(a-2)b$ and $c'=\operatorname{sgn}(a)c$.
Then
\[
                       t^2-Atc'+(c')^2=1.
\]
We have $|c'|\ge|a-2|\ge6$, so $t\ne0$. The two terms $t,c'$
have the same sign: opposite signs would make the left-hand side
at least $t^2+(c')^2>1$. Thus $T=|t|$ and $C=|c'|$ are positive
integer solutions of $T^2-ATC+C^2=1$.

Define $U_0=0$, $U_1=1$ and $U_{j+1}=AU_j-U_{j-1}$.
Every positive solution consists of adjacent terms of this sequence.
Indeed, order its coordinates as $M\ge m>0$. The integer
\[
                     Am-M=\frac{m^2-1}{M}
\]
lies in $[0,m)$, and replacing $(M,m)$ by $(m,Am-M)$ preserves
the equation. Repeated descent ends at $(1,0)$. Reversing the
steps gives the unordered pairs $\{U_{j+1},U_j\}$. Consequently
$C=U_j$ for some $j\ge1$ and $T\ge U_{j-1}$.

Put $d=|a-2|$. Since $d\mid C$, reduction of the recurrence
modulo $d$ gives $d\mid j$: if $a>0$, then $A\equiv2\pmod d$
and $U_j\equiv j\pmod d$; if $a<0$, then $A\equiv-2\pmod d$
and $U_j\equiv(-1)^{j-1}j\pmod d$. Hence $j\ge d\ge A-2$.
The recurrence also gives $U_r>(A-1)^{r-1}$ for $r\ge2$.
This follows by induction from $U_2=A>A-1$ and
$U_{r+1}>(A-1)U_r$, since $U_r>U_{r-1}$.
As $j\ge6$, we obtain
\[
 \log T>(j-2)\log(A-1)\ge(A-4)\log7>A-3
 \qquad(A\ge8).
\]
The last inequality holds at $A=8$ and its difference increases
with $A$. This proves the assertion.
\end{proof}

\subsection{An exact height identity}

\begin{lemma}\label{lem:u29-product}
For any nonzero odd integers $n_1,\ldots,n_r$,
\[
 h\left(\prod_{i=1}^r x([n_i]P)\right)
       =\sum_{i=1}^r h(x([n_i]P)).
\]
This identity requires only that the displayed $P=2R_0$ have infinite
order; it does not require the contact-height or prime-pattern hypotheses.
\end{lemma}
\begin{proof}
First consider a finite place $v$ with $v(D)=0$.  The formal subgroup
$E_1(K_v)$ consists of the identity and points with $v(x)<0$.
It is an open subgroup, so $E(K_v)/E_1(K_v)$ is finite.  These facts
hold for the displayed integral Weierstrass equation, including the
place above two.  The image of $T_0$ has order two, since
$T_0\notin E_1(K_v)$, and
\[
                    x(Q+T_0)=-D/x(Q).
\]
If an odd multiple of $P$ belongs to $E_1(K_v)$, the image of $P$
in the quotient has odd order.  No multiple of that image is the
nonzero element represented by $T_0$.  Hence no odd multiple of $P$
has positive $x$-valuation.  If no odd multiple belongs to $E_1(K_v)$,
all odd multiples have nonnegative $x$-valuation.  Thus the numbers
$v(x([n_i]P))$ have a common weak sign.

The only remaining finite place is the prime $(3-\sqrt2)$ above seven.
Here $v(D)=2$, and the discriminant has valuation six, so the displayed
integral model is minimal.  Its reduction is $y^2=x^3$.
The displayed $P$ reduces to $(1,6)$, a nonsingular point.  The subgroup
of points with nonsingular reduction therefore contains every multiple
of $P$.  A point with positive $x$-valuation would have positive
$y$-valuation and reduce to the singular point $(0,0)$.  Consequently
every multiple of $P$ has $x$-valuation at most zero at this place.

At either real embedding every $[n_i]P$ is a nonzero double.  The
duplication identity
\[
 x(2Q)=\left(\frac{x(Q)^2+D}{2y(Q)}\right)^2
\]
makes its abscissa positive.  The curve equation and its nonzero
ordinate imply $x([n_i]P)>d>1$ in each embedding.  Each real local
height therefore adds under multiplication.  At a finite place the
common weak sign similarly gives
\[
 \max\{0,-\textstyle\sum_i v(x([n_i]P))\}
       =\sum_i\max\{0,-v(x([n_i]P))\}.
\]
Summing with the normalized local weights proves the identity.
\end{proof}

\begin{corollary}\label{cor:u29-scales}
Let $Q_1,\ldots,Q_4$ be odd multiples of $P$ and write
\[
 C=\prod_{i=1}^4x(Q_i)=C_0+C_1\sqrt2.
\]
The three numbers $A_0=1$, $A_1=C_0^2$, and
$A_2=(1+C_1^2)^2$ are nonzero positive rational squares.
If the three contact detectors at $A_j\lambda$ imply
$M_0(A_j\lambda)^k\mid N\ne0$, then
\[
                         h(x(Q_i))<7H|N|^{19/k}
                         \qquad(1\le i\le4).
\]
Here $M_0$ is the odd-contact radical for
$\Theta=\{0,-1/2,-1,1/2399,-32/407\}$.
\end{corollary}
\begin{proof}
The product is positive in both embeddings, so $C_0>0$.
The fractional linear map $t\mapsto2+1/t$ identifies the five
contacts with the original five, outside the exceptional primes.
Consequently $h(t)\le(H+\log3)M_0(t)^{19}$.  Put
$T=(H+\log3)|N|^{19/k}$.  The unscaled detector and each scaled
detector give $h(A_j)\le2T$.  Hence $h(C_0)\le T$.
For $C_1=u/v$ in lowest terms with $v>0$,
$h(1+C_1^2)=\log(u^2+v^2)\ge2h(C_1)$, so $h(C_1)\le T/2$.
The elementary height inequality in the fixed basis gives
\[
 h(C)\le\tfrac32T+\tfrac32\log2<7H|N|^{19/k}.
\]
Lemma~\ref{lem:u29-product} bounds every summand by $h(C)$.
\end{proof}

\section{The exact numerical height constant}
\label{sec:exact-height-constant}

The contact-height input uses a fixed, explicitly specified integer.
Its exact evaluation improves the numerical bounds in both construction
methods. Put
\[
 B_*=\prod_{p\le76176}p,\qquad
 H_{19}=2^{110}210^{1800}B_*^{19}+128.
\]
This is the constant of the source's exponent-nineteen height theorem
\cite[Theorem \nolinkurl{thm:quartic-height19}]{RepositoryV9}.

\begin{lemma}\label{lem:exact-height-constant}
The exact inequalities
\[
                 2^{2093447}<H_{19}<2^{2093448}
\]
hold. In particular one may use $h=2093448$ in every argument
requiring $H_{19}<2^h$.
\end{lemma}
\begin{proof}
This finite integer calculation has the following reproducible
certificate. Enumerate primes through $76176$ by marking, for
each unmarked $p\le\lfloor\sqrt{76176}\rfloor$, all multiples
starting at $p^2$. The remaining integers at least two are
exactly the primes: every composite has a prime divisor not
exceeding its square root, and no prime is marked.
An independent enumeration tests every integer by trial division
through its integer square root. Both lists have $7499$ entries
and agree entry by entry.

Their exact product has $109445$ binary digits. Computing the
displayed powers and product as integers gives $2093448$ binary
digits for $H_{19}$. The comparison with the next smaller power
of two is strict as well. The supplied checker performs both
prime enumerations and both integer comparisons; the complete
prime list and the unsigned binary encoding of $H_{19}$ are
included. No floating-point logarithm enters these comparisons.
\end{proof}

This lemma evaluates a constant. It does not prove the
contact-height inequality itself. The latter is still the
arithmetic input stated in Section~\ref{sec:u29}, now used with
this exact value and its sharper certified bound.


\section{Simultaneous parity and divisibility conditions}
\label{sec:parity-batching}

Put $K=\mathbb Q(\rho)$, where $\rho^2=2$, and consider
\[
 d=3-\rho,\qquad D=d^2=11-6\rho,\qquad
 E:\ y^2=x^3-Dx,
 \qquad
 P=\left(\frac{57-28\rho}{8},
          \frac{-660+417\rho}{32}\right).
\]
We write $v_p$ for the rational valuation with $v_p(p)=1$.
The valuation $\nu_2$ on $K$ above $2$ is normalized by
$\nu_2(\rho)=1$, so $\nu_2(2)=2$.
Let $R=\mathbb Z_{(2)}\cap\mathbb Z_{(3)}$, and set
\[
                         N(a,b)=a^2+ab-b^2.
\]

\begin{lemma}\label{lem:four-semilocal-inputs}
For rational $a,b,c,e$,
\[
 N(a,b)+6N(c,e)\in R
 \quad\Longleftrightarrow\quad a,b,c,e\in R.
\]
\end{lemma}
\begin{proof}
The binary form $N$ has no nontrivial zero over either $\mathbb F_2$
or $\mathbb F_3$. Consequently, for $p=2,3$ and $(a,b)\ne(0,0)$,
\[
                    v_p(N(a,b))=2\min(v_p(a),v_p(b)).
\]
Indeed, divide both inputs by the smaller power of $p$ and reduce
modulo $p$. The two summands $N(a,b)$ and $6N(c,e)$ have valuations
of opposite parity when both are nonzero. They therefore cannot
cancel. Their sum is $p$-integral if and only if
\[
 2\min(v_p(a),v_p(b))\ge0,
 \qquad 1+2\min(v_p(c),v_p(e))\ge0.
\]
Since these minima are integers, these inequalities say that all four
inputs are $p$-integral. If either pair is zero, the same conclusion
follows by omitting the corresponding zero term.
\end{proof}

\begin{lemma}\label{lem:semilocal-three-witnesses}
For $t\in\mathbb Q$,
\[
 t\in R\quad\Longleftrightarrow\quad
 \exists u,v,w\in\mathbb Q:\
                u^2+6v^2-2w^2-3t^2+2=0.
\]
\end{lemma}
\begin{proof}
This is the three-witness semilocal criterion used by the source
construction; we recall its arithmetic justification.
The Hilbert-symbol formulas show that $(2,3)_{\mathbb Q}$ is
ramified precisely at $2$ and $3$: outside $2,3$ both parameters
are units and the algebra splits; at $3$ the symbol is the
Legendre symbol $(2/3)=-1$, and reciprocity gives the symbol
$-1$ at $2$ since the real place is split. Its quadratic-field
embedding criterion \cite[Proposition 14.6.7]{Voight2021} says that a
nonzero rational number $\Delta$ is the square of a trace-zero
element precisely when $\Delta$ is nonsquare in both completions.
For a rational square $\Delta\ne0$, such an element is impossible
in a division algebra: $(q-\sqrt\Delta)(q+\sqrt\Delta)=0$ would
make $q$ scalar. For nonsquare $\Delta$, this is the usual local
embedding criterion for quadratic fields in quaternion algebras.

For $t\ne0$, take $\Delta=1-2/(3t^2)$, which is nonzero because
$2/3$ is not a rational square. The trace-zero square form is
$2x^2+3y^2-6z^2$. At $2$, if $v_2(t)>0$, then $v_2(\Delta)$ is
odd; if $v_2(t)=0$, then $\Delta\equiv3\pmod8$; and if
$v_2(t)<0$, then $\Delta\in1+8\mathbb Z_2$ is a square.
At $3$, a nonnegative valuation of $t$ gives odd negative valuation
of $\Delta$, whereas a negative valuation gives a square in
$1+3\mathbb Z_3$. Hence the embedding criterion holds exactly
when $t\in R$. Multiplying
$2x^2+3y^2-6z^2=\Delta$ by $3t^2$, and putting
$v=tx$, $u=3ty$, $w=3tz$, gives the displayed equation.
The case $t=0$ has witnesses $(u,v,w)=(0,0,1)$.
\end{proof}

\begin{lemma}\label{lem:exact-dyadic-index}
For every nonzero integer $n$,
\[
 \nu_2(x([n]P))=-6-4v_2(n).
\]
At the place above $3$, every odd multiple $[n]P$ has unit abscissa.
\end{lemma}
\begin{proof}
The displayed coordinates give
$\nu_2(x(P))=-6$, $\nu_2(y(P))=-9$, and therefore
$\nu_2(z(P))=3$ for the formal parameter $z=-x/y$.
The equation is integral at $2$. Its formal logarithm is an isometry
on the subgroup of depth strictly greater than
$\nu_2(2)/(2-1)=2$. No minimal-model assertion is needed for this
integral formal group \cite[Theorem IV.6.4]{Silverman2009}.
Its nonzero logarithm also proves that $P$ has infinite order.
Since the logarithm is a homomorphism,
\[
 \nu_2(z([n]P))=3+2v_2(n).
\]
The Laurent expansion $x=z^{-2}(1+O(z))$ then proves the formula.

At $3$ the equation has good reduction: its discriminant is
$64D^3$, a unit. The residue field is $\mathbb F_9$, and direct
reduction gives $\overline P=(\overline\rho,0)$.
This point has order two and nonzero abscissa. Every odd multiple
therefore has the same reduction and unit abscissa.
\end{proof}

\begin{theorem}\label{thm:parity-product}
Let $n_1,\ldots,n_r$ be nonzero integers and write
\[
        C=\prod_{i=1}^{r}x([n_i]P)=C_0+C_1\rho.
\]
Then
\[
 n_1,\ldots,n_r\text{ are all odd}
 \quad\Longleftrightarrow\quad
                2^{3r}C_0,\ 2^{3r}C_1\in R.
\]
\end{theorem}
\begin{proof}
The preceding lemma gives
\[
              \nu_2(2^{3r}C)=-4\sum_{i=1}^{r}v_2(n_i).
\]
Each summand on the right is nonnegative. Thus this expression is
nonnegative exactly when every index is odd.
The local integer rings at $2$ and $3$ are respectively
$\mathbb Z_2[\rho]$ and $\mathbb Z_3[\rho]$.
At $2$ this can also be seen directly from
\[
 \nu_2(a+b\rho)=\min(2v_2(a),2v_2(b)+1),
\]
whose two entries have opposite parity. At $3$, the integral basis
$(1,\rho)$ remains linearly independent modulo $3$.
Therefore integrality in either local ring is equivalent to
integrality of both rational coefficients. If the indices are odd,
the preceding lemma makes $C$ a $3$-adic unit, so the additional
$3$-integrality conditions impose no restriction. The converse
already follows from the place above $2$.
\end{proof}



\section{Prescribing primary divisibility with the existing anchor}
\label{sec:primary-depth-anchor}

This section concerns rational points on one fixed curve.  Its local
statements do not use the prime-pattern or height inputs of the
Diophantine construction.  Those arithmetic inputs enter only when
the construction supplies the specified cyclic point domain and its
integer-input witnesses.

Put
\[
 K=\mathbb Q(\sqrt2),\qquad D=11-6\sqrt2,\qquad
 E:y^2=x^3-Dx,
\]
and
\[
 P=\left(\frac{57}{8}-\frac72\sqrt2,
          -\frac{165}{8}+\frac{417}{32}\sqrt2\right).
\]
At the place above two, normalize the valuation $\nu_2$ by
$\nu_2(\sqrt2)=1$, so $\nu_2(q)=2v_2(q)$ for $q\in\mathbb Q^*$.
Write
\[
 R=\mathbb Z_{(2)}\cap\mathbb Z_{(3)},\qquad
 N(a,b)=a^2+ab-b^2.
\]
All points whose reciprocal abscissas occur below are finite and
nonidentity.  All cyclic indices are nonzero integers.

\begin{lemma}\label{pdanchor:local-coefficients}
For every nonzero integer $n$, write
$f([n]P)=1/x([n]P)=a_n+b_n\sqrt2$.  Then
\[
 \nu_2\bigl(x([n]P)\bigr)=-6-4v_2(n),\qquad
 v_2(a_n)=3+2v_2(n),\qquad
 v_2(b_n)\ge3+2v_2(n).
\]
At the place above three, an odd multiple of $P$ has unit
abscissa, whereas an even nonzero multiple is formal and its
reciprocal abscissa has both rational coefficients in $\mathbb Z_3$.
For $p=2,3$ and rational $a,b$, not both zero, one has
\[
             v_p\bigl(N(a,b)\bigr)=2\min(v_p(a),v_p(b)).
\]
\end{lemma}
\begin{proof}
The equation for $E$ is integral at two.  The displayed coordinates
give $\nu_2(x(P))=-6$ and $\nu_2(y(P))=-9$.  Thus the formal
parameter $z=-x/y$ has $\nu_2(z(P))=3$.  The formal logarithm
is an isometry when $\nu_2(z)>2$, since the absolute ramification
index is two.  Multiplication by $n$ therefore gives
\[
 \nu_2(z([n]P))=3+2v_2(n),\qquad
 \nu_2(x([n]P))=-2\nu_2(z([n]P)).
\]
The logarithm also shows that these multiples are nonidentity.
For rational $a,b$, the two quantities $2v_2(a)$ and
$2v_2(b)+1$ have different parity.  Hence
\[
 \nu_2(a+b\sqrt2)=\min\{2v_2(a),2v_2(b)+1\}.
\]
Applied to $f([n]P)$, whose valuation is the even integer
$6+4v_2(n)$, this proves the two coefficient assertions.

At three the curve has good reduction and
$\overline P=(\sqrt2,0)$ has order two.  Its abscissa is a unit.
This proves the assertions for odd and even multiples.  The
extension $\mathbb Q_3(\sqrt2)/\mathbb Q_3$ is unramified,
with integral basis $(1,\sqrt2)$, so integrality of an element
implies integrality of both displayed coefficients.

Finally, $X^2+XY-Y^2$ has no nonzero zero over either
$\mathbb F_2$ or $\mathbb F_3$.  After factoring the common
minimum power of $p$ from $a$ and $b$, reduction modulo $p$
therefore gives a unit.  This proves the norm valuation formula.
\end{proof}

\begin{lemma}\label{pdanchor:seven-depths}
Let $1\le s\le7$ and $r\ge1$.  For points $U_i=[u_i]P$,
write $f(U_i)=a_i+b_i\sqrt2$ and put
\[
 T_r=2^{-(4r+5)}\sum_{i=1}^{s}a_i^2.
\]
Let $m=\min_i v_2(u_i)$, and let $k$ be the number of indices
attaining this minimum.  Then
\[
                 v_2(T_r)=1+4(m-r)+v_2(k).
\]
Consequently $T_r$ is $2$-integral if and only if
$2^r\mid u_i$ for every $i$.  Whenever $T_r$ is not integral,
its valuation is congruent to $1$, $2$, or $3$ modulo four.
\end{lemma}
\begin{proof}
By Lemma~\ref{pdanchor:local-coefficients}, the numbers
$w_i=a_i/2^{3+2m}$ are $2$-integral.  Exactly $k$ of them
are units, and every other $w_i$ is divisible by four.
A $2$-adic unit square is congruent to one modulo eight, so
\[
                   \sum_iw_i^2\equiv k\pmod8.
\]
Since $1\le k\le7$, one has $v_2(k)\le2$, and the displayed
congruence determines the valuation of this sum exactly.
Rescaling gives
\[
 v_2(T_r)=6+4m-(4r+5)+v_2(k)
         =1+4(m-r)+v_2(k).
\]
This is positive for $m\ge r$ and negative for $m<r$.
Its residue modulo four belongs to $\{1,2,3\}$.
The sum is nonzero, also directly because its rational summands
are squares and every $a_i$ is nonzero.
\end{proof}

The bound of seven is necessary for this particular unweighted
test.  For eight identical points $U_i=P$ and $r=1$, one has
$T_1=a_1^2/64$ and $v_2(T_1)=0$, although all eight indices
are odd.

\begin{theorem}\label{pdanchor:merged-depth}
Let $Q_j=[n_j]P$, $1\le j\le4$, be four nonzero cyclic
points, and write
\[
 C=\prod_{j=1}^4x(Q_j)=C_0+C_1\sqrt2,\qquad
 B=2^{24}N(C_0,C_1).
\]
Let $S=[a]P$, with $a\ne1$, and write
$f(S-P)=f_0+f_1\sqrt2$.  Fix a nonnegative integer $q$ and
put $C_*=3^q$, $S_0=C_*^2$, and
\[
                         A=\frac{N(f_0,f_1)}{S_0}.
\]
With $1\le s\le7$ primary points and $T_r$ as in
Lemma~\ref{pdanchor:seven-depths}, the condition
\[
                            A+B+T_r\in R
\]
is equivalent, on the stated cyclic point domain, to all three
of the following conditions:
\[
 \frac{f_0}{C_*},\frac{f_1}{C_*}\in R,
 \qquad n_j\text{ is odd for every }j,
 \qquad 2^r\mid u_i\text{ for every }i.
\]
\end{theorem}
\begin{proof}
Set $M=\sum_jv_2(n_j)$.  Lemma~\ref{pdanchor:local-coefficients}
gives
\[
                  \nu_2(2^{12}C)=-4M.
\]
The valuation is even.  Therefore
$v_2(2^{12}C_0)=-2M$ and
$v_2(2^{12}C_1)\ge-2M$; the norm valuation formula yields
\[
                             v_2(B)=-4M.
\]
The point $S-P=[a-1]P$ has nonzero integer index.  Both
coefficients of its reciprocal abscissa are $2$-integral, and
$C_*$ is a $2$-adic unit.  Thus $A$ is always $2$-integral.

If $B$ is not $2$-integral, its valuation is a negative
multiple of four.  If $T_r$ is not $2$-integral, its valuation
is not a multiple of four.  Two such negative valuations are
different, so their minimum is the valuation of their sum.
If only one of $B,T_r$ is nonintegral, its negative valuation
likewise cannot be canceled by the other.  Adding the integral
term $A$ cannot change either conclusion.  Hence
\[
 A+B+T_r\in\mathbb Z_2
 \quad\Longleftrightarrow\quad
 B,T_r\in\mathbb Z_2.
\]
The right side says exactly that all $n_j$ are odd and all
$u_i$ are divisible by $2^r$.

Assume these conclusions.  At three, each root point reduces
to $(\sqrt2,0)$, so $C_0,C_1$ and $B$ are $3$-integral.
Since $r\ge1$, every primary point is an even multiple and
reduces to the identity.  Its reciprocal abscissa has integral
rational coefficients, and the denominator of $T_r$ is a
power of two.  Thus $T_r$ is also $3$-integral.  Subtraction
now gives
\[
 A+B+T_r\in\mathbb Z_3
       \quad\Longleftrightarrow\quad A\in\mathbb Z_3.
\]
The norm valuation formula identifies the last condition with
the simultaneous $3$-integrality of $f_0/C_*$ and $f_1/C_*$.
Their $2$-integrality was automatic.  Conversely, the three
stated conditions put each of $A,B,T_r$ in $R$.
\end{proof}

\paragraph{Quadratic implementation.}
The theorem uses the single physical-anchor value
\[
 G=N(f_0,f_1)+S_0\left(2^{24}N(C_0,C_1)
                     +2^{-(4r+5)}\sum_{i=1}^{s}a_i^2\right).
\]
This is a quadratic expression in coordinates already present.
The existing ternary semilocal representation imposes
$G/S_0\in R$ with its existing three witnesses.  No additional
coordinate is needed.  For $r=1$, an invertible rational
rescaling $B_j=2^{12}C_j$ and $A_i=a_i/32$ writes the
physical-value equation as
\[
 N(f_0,f_1)+S_0\left(N(B_0,B_1)+2\sum_iA_i^2\right)-G=0.
\]
The coordinate rescalings preserve the degree after fixed
denominators are cleared.

\begin{corollary}\label{pdanchor:five-inverse-pairs}
In the construction with five primary points $U_0,\ldots,U_4$,
take $r=1$ and $Q_0=P$, with $Q_1,\ldots,Q_4$ the four root
points.  On the strengthened point domain of
Theorem~\ref{pdanchor:merged-depth}, the two rational coordinates
for each inverse $(x(U_i)-x(Q_i))^{-1}$ can be deleted.
If these ten coordinates occur only in their inverse equations,
their deletion changes neither the solutions of this strengthened
system nor its input projection.
\end{corollary}
\begin{proof}
The theorem makes each $U_i$ an even multiple of $P$ and
each $Q_i$ an odd multiple.  Equality of their abscissas would
imply $U_i=Q_i$ or $U_i=-Q_i$ on the short Weierstrass curve.
The infinite order of $P$ would then identify an even integer
with an odd one, which is impossible.  Each difference is
therefore nonzero.  The omitted inverse exists uniquely in
$K$, giving its two rational coordinates and uniquely extending
every solution of the reduced system.  In particular, the
retained slope equations have the same intended values.
\end{proof}

The comparison with the earlier construction requires a separate
completeness argument.  Its integer-input witnesses use primary
indices $u_i=hn_i$ with $64\mid h$; all five are already even.
They therefore satisfy the strengthened domain and the new merged
anchor.  Under the inherited arithmetic hypotheses this preserves
completeness at every integer input, while every solution of the
strengthened system remains a valid solution of the earlier one.
It is not asserted that every auxiliary solution of the earlier
system has even primary indices.  The ten-coordinate saving uses
this proved source restriction, together with the independently
checked occurrence condition on the deleted inverse coordinates.


\section{Exact three-adic depth of the anchor}
\label{sec:three-adic-anchor-depth}

The coefficient size of the physical anchor can be reduced by using its
exact local depth.  The following calculation concerns only the fixed
point and its integral multiples; it is independent of the prime-pattern
and global height hypotheses.

Put
\[
 K=\mathbb Q(\sqrt2),\qquad D=11-6\sqrt2,\qquad
 E:y^2=x^3-Dx,
\]
and
\[
 P=\left(\frac{57}{8}-\frac72\sqrt2,
          -\frac{165}{8}+\frac{417}{32}\sqrt2\right).
\]
The prime three is inert in $K$.  Let $v$ be the valuation on
$K_3=\mathbb Q_3(\sqrt2)$ normalized by $v(3)=1$.
As usual, the valuation of zero is $+\infty$.

\begin{proposition}\label{anchor3:exact-depth}
For every nonzero integer $c$, the point $[c]P$ is finite and has
nonzero abscissa.  Write
\[
             f([c]P)=\frac1{x([c]P)}=a_c+b_c\sqrt2.
\]
Then
\[
 \min\{v_3(a_c),v_3(b_c)\}=
 \begin{cases}
 0,&c\text{ odd},\\
 2+2v_3(c),&c\text{ even}.
 \end{cases}
\]
In the even case the formal parameter $z=-x/y$ satisfies
\[
                  v(z([c]P))=1+v_3(c).
\]
\end{proposition}
\begin{proof}
The displayed equation has good reduction at three: its discriminant is
$64D^3$, which is a unit.  Direct reduction gives
\[
             \overline P=(\sqrt2,0)\in E(\mathbb F_9).
\]
This is a nonidentity point of order two, with unit abscissa.  Hence
every odd multiple of $P$ is finite with unit abscissa.

The ordinate of $P$ has valuation one, since
\[
 y(P)=\frac{3(-220+139\sqrt2)}{32}
\]
and its parenthesized factor has nonzero residue in $\mathbb F_9$.
The numerator of the tangent slope is a unit: modulo three,
$3x(P)^2-D=-D\ne0$.  Thus
\[
 \lambda=\frac{3x(P)^2-D}{2y(P)}
 \quad\hbox{satisfies}\quad v(\lambda)=-1.
\]
The duplication formula now gives
\[
 \begin{split}
 v(x([2]P))&=v(\lambda^2-2x(P))=-2,\\
 v(y([2]P))&=v\bigl(\lambda(x(P)-x([2]P))-y(P)\bigr)=-3.
 \end{split}
\]
There is no cancellation: the summands in the first difference have
valuations $-2,0$, and those in the second have valuations $-3,1$.
Consequently $v(z([2]P))=1$.

For completeness, the needed formal-group valuation is elementary here.
The formal group of the integral Weierstrass equation has integral
coefficients, and negation sends $z$ to $-z$.  Its multiplication series
therefore have the form
\[
                  [m](Z)=mZ+Z^3g_m(Z^2),
                  \qquad g_m(T)\in\mathcal O_{K_3}[[T]].
\]
If $v(z)\ge1$ and $3\nmid m$, the linear term has strictly smaller
valuation than every higher term, so $v([m]z)=v(z)$.
For $m=3$, the linear term has valuation $v(z)+1$, whereas all higher
terms have valuation at least $3v(z)>v(z)+1$.  Thus
\[
                    v([3]z)=v(z)+1.
\]
Write $c=2\cdot3^ru$, where $3\nmid u$.  Repeated application to
$z([2]P)$ yields $v(z([c]P))=1+r$.  In particular, this parameter
is nonzero, so $[c]P$ is not the identity.  On the formal group,
$x=z^{-2}$ times an integral power series with constant term one.
It follows that
\[
                    v(f([c]P))=2+2r.
\]

Finally, $1,\sqrt2$ form an integral basis for the unramified quadratic
extension $K_3/\mathbb Q_3$, with linearly independent residues over
$\mathbb F_3$.  For any rational $a,b$, not both zero, this gives the
exact equality
\[
                  v(a+b\sqrt2)=\min\{v_3(a),v_3(b)\}.
\]
Applying it to the reciprocal abscissas proves both cases.
\end{proof}

\begin{corollary}\label{anchor3:integer-cutoff}
Let $N\ge1$ and $c\ne0$.  Then
\[
 \frac{a_c}{3^N},\frac{b_c}{3^N}\in\mathbb Z_3
 \quad\Longleftrightarrow\quad
 c\text{ is even and }v_3(c)\ge\left\lceil\frac N2\right\rceil-1.
\]
For a prescribed integer $r\ge0$, the smallest positive $N$ forcing
$v_3(c)\ge r$ is $N=2r+1$.  At this value the anchor implies
\[
                             |c|\ge2\cdot3^r.
\]
\end{corollary}
\begin{proof}
The displayed divisibility condition says that the minimum valuation in
Proposition~\ref{anchor3:exact-depth} is at least $N$.  Since $N>0$,
this excludes the odd case.  In the even case it is exactly
$2+2v_3(c)\ge N$.  The stated rounding follows.  When $r\ge1$,
the value $N=2r$ still admits $c=2\cdot3^{r-1}$, so $2r+1$ is
minimal; for $r=0$, it is the smallest positive value by definition.
The final bound follows from the simultaneous divisibility by two and
$3^r$.
\end{proof}

\paragraph{Application to the physical anchor.}
In the merged-anchor theorem, replace its arbitrary constant $C_*$ by
$3^N$ and its coefficient $S_0=C_*^2$ by $3^{2N}$.
The two-adic argument is unchanged because three is a two-adic unit.
Once the root parity and primary divisibility conditions have been
extracted, the three-adic argument recovers exactly the two conditions
in Corollary~\ref{anchor3:integer-cutoff}, with $c=a-1$.
The existing point domain requires $C_S=S-P\ne O$, so $c\ne0$;
this condition must be retained.  If one instead extended $f$ by
$f(O)=0$, the anchor alone would not exclude $c=0$.

Suppose the inherited global argument needs
\[
                     |c|>(80H)^{10},\qquad H<2^h.
\]
It suffices to choose an integer $r$ satisfying
\[
                       2\cdot3^r>(80\cdot2^h)^{10}
\]
and to take
\[
                   N=2r+1,\qquad C_*=3^{2r+1},
                   \qquad S_0=3^{4r+2}.
\]
These substitutions change only fixed integer coefficients.  They add
no coordinate or equation and do not change the polynomial degree.
Completeness is retained by requiring the already freely chosen source
index $c$ to be divisible by $2\cdot3^r$ in the same fixed local
congruence conditions.  This is a fixed refinement of the construction,
not a uniformity assertion with an increasing prime-pattern modulus.

\begin{corollary}\label{anchor3:numerical-constant}
For $h=2093448$, the smallest integer $r$ satisfying
$2\cdot3^r>(80\cdot2^h)^{10}$ is
\[
                         r=13208226.
\]
The corresponding smallest anchor exponent is
\[
 N=26416453,\qquad C_*=3^{26416453},\qquad
 S_0=3^{52832906}.
\]
In particular, a coefficient-height formula $64S_0^2$ specializes to
\[
                              64\cdot3^{105665812}.
\]
The minimality assertion concerns the stated conservative bound
$H<2^{2093448}$ and this exact anchor criterion.
\end{corollary}
\begin{proof}
The threshold for $r$ is equivalently
\[
                     3^r>5^{10}2^{20934519}.
\]
The following four strict integer comparisons are verified by exact
integer exponentiation in the accompanying checker:
\[
 \begin{array}{ll}
 3^{31867}>2^{50508},&
 3^{15288}>5^{10}2^{24207},\\
 3^{15601}<2^{24727},&
 3^{9779}<5^{10}2^{15477}.
 \end{array}
\]
The integer identities
\[
 \begin{split}
 13208226&=414\cdot31867+15288,\\
 20934519&=414\cdot50508+24207
 \end{split}
\]
give the required lower comparison.  Similarly,
\[
 \begin{split}
 13208225&=846\cdot15601+9779,\\
 20934519&=846\cdot24727+15477
 \end{split}
\]
give $3^{13208225}<5^{10}2^{20934519}$ and hence minimality.
The remaining values follow from $N=2r+1$ and $S_0=3^{2N}$.
\end{proof}

The exact checker accompanying this calculation verifies the source
point, its duplication formulas and valuations, and finite examples of
both parity cases using rational arithmetic in $\mathbb Q(\sqrt2)$.
The statement for all indices follows from the formal-group proof,
not from extrapolating those examples.


\section{A scalar normalization without squareclass corrections}
\label{sec:scalar-correction-free}

Put
\[
                         \beta(X)=216X^4+450459
\]
and retain the five-contact polynomial
\[
 r(T)=T(T-1)(T-2401)(32T+343)(2T-1).
\]
Define
\[
 r_X=r(\beta(X)),\qquad
 \delta_X=1+(17r_X)^2,\qquad
 R_X=\frac{102^2r_X}{\delta_X^2}.
\]
The scaling factor is an integer.  In particular, this normalization
introduces no fixed denominator prime.

\begin{lemma}\label{sn:integral-input}
For $X\in\mathbb Q$ one has
\[
                 \beta(X)\in\mathbb Z
                 \quad\Longleftrightarrow\quad X\in\mathbb Z.
\]
For every real $X$, $\beta(X)\ge450459>2401$, so $r_X>0$
and $\delta_X>0$.  If $p>76176$ and $v_p(\beta(X))<0$,
then
\[
 v_p(\beta(X))=4v_p(X),\qquad
 v_p(r_X)=5v_p(\beta(X)),\qquad
 v_p(R_X)=-3v_p(r_X).
\]
In particular, all three valuations are even.
\end{lemma}
\begin{proof}
If $X$ is an integer, so is $\beta(X)$.  Conversely, write
$X=a/b$ in lowest terms with $b>0$.  Integrality of $\beta(X)$
gives $b^4\mid216=2^3\cdot3^3$.  Neither prime exponent
admits a nontrivial fourth power, so $b=1$.

The real inequalities follow from the fourth power and the
five positive factors of $r(T)$ when $T>2401$.  At a prime
in the stated range, a negative valuation of $\beta$ forces
$v_p(X)<0$.  The unique dominant terms in $\beta$ and
$r(\beta)$ are respectively $216X^4$ and $64\beta^5$.
This gives the first two valuation formulas.  In this case
$v_p(r_X)<0$, so the dominant term of $\delta_X$ is
$17^2r_X^2$, and $v_p(\delta_X)=2v_p(r_X)$.  The final
formula follows.  The displayed multiple of four is even.
\end{proof}

\begin{lemma}\label{sn:local-squares}
If $X\in\mathbb Z$, then $r_X$ is a nonzero square in
$\mathbb Q_2$, $\mathbb Q_3$, and $\mathbb Q_{17}$.
More precisely,
\[
 v_2(r_X)=2,\quad r_X/4\equiv1\pmod8,
 \qquad
 v_3(r_X)=2,\quad r_X/9\equiv1\pmod3,
\]
and at seventeen either $v_{17}(r_X)=0$ with square-unit
residue, or
\[
                    v_{17}(r_X)=2,\qquad
                    r_X/17^2\equiv9\pmod{17}.
\]
\end{lemma}
\begin{proof}
Since $216$ is divisible by eight and by twenty-seven,
\[
                    \beta\equiv3\pmod8,
                    \qquad\beta\equiv18\pmod{27}.
\]
For the two even factors of $r(\beta)$ one has
\[
 v_2(\beta-1)=v_2(\beta-2401)=1,
 \qquad
 \frac{\beta-1}{2}-\frac{\beta-2401}{2}=1200\equiv0\pmod8.
\]
Their odd parts therefore have square product congruent to one
modulo eight.  The remaining odd factors have product
$3\cdot7\cdot5\equiv1\pmod8$.  This proves the assertions
at two.  At three, only the factor $\beta$ is divisible by
three; it has valuation two.  Reduction after division by nine
gives
\[
 \frac{r(\beta)}9\equiv
       2\cdot2\cdot2\cdot1\cdot2\equiv1\pmod3.
\]

At seventeen the fourth powers are exactly $0,1,4,13,16$.
Since $216\equiv12$ and $450459\equiv10\pmod{17}$, all
possibilities are given by the table
\[
\begin{array}{c|rrrrr}
 X^4&0&1&4&13&16\\\hline
 \beta&10&5&7&13&15\\
 r(\beta)&0&15&2&13&2.
\end{array}
\]
The nonzero residues are squares: $13=8^2$, $15=7^2$, and
$2=6^2$ modulo seventeen.

In the remaining case $17\mid X$.  The exact identity
\[
 32\cdot450459+343=17^2\cdot49879,
 \qquad 49879\equiv1\pmod{17},
\]
and $v_{17}(X^4)\ge4$ show that the factor
$32\beta+343=6912X^4+17^2\cdot49879$ has valuation exactly
two and unit residue one.  The other four factors are units,
with product
$10\cdot9\cdot6\cdot2\equiv9\pmod{17}$.
This proves the last assertion.  The usual local square criteria
at two and at odd primes now prove all three square statements.
\end{proof}

\begin{proposition}\label{sn:prepared-constant}
For every integer $X$, the number $R_X$ is prepared: it is
positive and less than one, its reduced denominator is a square,
it is a square in $\mathbb Q_2,\mathbb Q_3,\mathbb Q_{17}$,
and
\[
 v_2(R_X)=4,\qquad v_3(R_X)=4,\qquad
 v_{17}(R_X)\in\{2,4\}.
\]
No squareclass correction is required.
\end{proposition}
\begin{proof}
The contact value $r_X$ is a positive integer.  One has
$\gcd(r_X,\delta_X)=1$, and the preceding lemma shows that
$\delta_X$ is a unit at two and three; it is a unit at
seventeen by its definition.  Thus
$\gcd(102^2r_X,\delta_X^2)=1$, and the reduced denominator
is exactly $\delta_X^2$.

The quotient $R_X/r_X=(102/\delta_X)^2$ is a rational square.
Its local valuations add two at each of $2,3,17$, proving
both the asserted squareclasses and the exact depths.  Finally,
since $r_X\ge1$,
\[
 0<R_X=
 \frac{102^2r_X}{(1+17^2r_X^2)^2}
 <\frac{102^2}{17^4r_X^3}
 \le\frac{10404}{83521}<1.
\]
\end{proof}

\subsection{The norm tree and its domain}

Fix a positive rational number $\rho$ and six rational tests
$t_0,\ldots,t_5$, with their assigned positive integer orders.
Start the scalar norm tree with the distinguished node
$1+t_0^2/\rho$.  Its other initial nodes are the assigned
normalized tests or pair norms.  Every later operation is
\[
                  N_\rho(u,v)=u^2+\rho v^2
\]
or the unchanged carrying of an unpaired node.  This is the
constant-left tree used by the scalar certificate; denote its
root by $W_\rho$.

\begin{lemma}\label{sn:positive-root}
Where its rational expressions are defined, $W_\rho>0$.
When all six tests tend to zero, $W_\rho$ tends to one.
For fixed nonzero rational $\rho$, at any finite set of primes
the tests can be required sufficiently small that
$W_\rho$ is a unit at every prime in that set.
\end{lemma}
\begin{proof}
The distinguished initial node is positive.  At each norm
operation its positive descendant contributes a positive
square, with a positive coefficient if it is the second
argument.  Its path therefore remains positive up to the
root.  When all tests vanish, this path starts at one and
all other nodes vanish, so its root is one.  For fixed
$\rho$, the tree is a polynomial in the normalized tests
with fixed rational coefficients.  It is continuous at zero
over $\mathbb R$ and over every $\mathbb Q_p$.
Its value at zero is one.  The inverse image of $\mathbb Z_p^*$
is therefore a neighborhood of zero, proving the final claim.
\end{proof}

There are two valid implementations.  One may replace the contact
parameter throughout the tree by $\rho=R_X$ and put
$b=R_XW_{R_X}^2$.  Alternatively, one may retain the original
tree with $\rho=r_X$ and normalize only its final pole input:
\[
                       b=R_XW_{r_X}^2.
\]
The second choice retains the original test denominators.  It
also has the quadratic scaling equations
\[
        \delta_X=1+17^2r_X^2,\qquad
        \delta_X V=102W_{r_X},\qquad b=r_XV^2,
\]
where the last expression can use the existing square-and-product
gates.  The first two displayed equations are quadratic.
The soundness and completeness proof for this substitution is
given below.

For rational inputs, both implementations have $r_X,R_X,\delta_X>0$
and $W_\rho>0$.  Thus the scalar contact, normalizer, norm-root,
and pole-input nonvanishing conditions follow on the stated
rational-expression domain.  A coordinate used solely to certify
one of these nonzero conditions can be removed after its occurrence
condition is checked.  This observation does not remove unrelated
denominator guards in the elliptic-point formulas.

\subsection{Soundness and simultaneous completeness}

\begin{proposition}\label{sn:contact-preservation}
At every finite scalar contact prime $p>76176$, the replacement
has the same positive odd contact valuation as the original
certificate.  Consequently it preserves its norm-tree extraction
of the assigned contact orders, for either implementation above.
\end{proposition}
\begin{proof}
At such a contact, $\nu=v_p(r_X)>0$ is odd.  The normalizer
satisfies $\delta_X\equiv1\pmod p$, and $102$ is a unit.
Hence
\[
               v_p(R_X)=v_p(r_X)=\nu,
 \qquad R_X/r_X\in\mathbb Z_p^{*2}.
\]
For either $\rho=r_X$ or $R_X$, the two terms in
$u^2+\rho v^2$ have valuations of opposite parity when both
are nonzero.  There is therefore no cancellation: its valuation
is the minimum of $2v_p(u)$ and $\nu+2v_p(v)$.
Every step of the existing constant-left extraction uses exactly
this identity, the assigned powers of the contact parameter in
the leaves, and the odd-pole exclusion for $b$.  All these
valuations are unchanged.  Lemma~\ref{sn:integral-input}
also shows that a contact with infinity has even valuation
outside the fixed exceptional set, so it creates no omitted
odd contact.  The inherited height and comparison arguments
therefore give $\beta\in\mathbb Z$, and
Lemma~\ref{sn:integral-input} gives $X\in\mathbb Z$.
\end{proof}

\begin{proposition}\label{sn:simultaneous-completeness}
Under the same arithmetic inputs as the source construction,
the scalar block with the input map and normalization above
can be satisfied at every integer input, together with the
remaining point, height, comparison, and size conditions.
The finite squareclass correction parameter can be omitted.
\end{proposition}
\begin{proof}
Fix an integer input.  The numbers $r_X,R_X,\delta_X$ are
now fixed, before any contact primes are chosen.  Proposition~
\ref{sn:prepared-constant} supplies the sign, local squareclasses,
square denominator, and small positive real size formerly obtained
by a correction.  Include the finite support of the numerator
and denominator of $R_X$, and that of $\rho$, in the fixed
local approximation conditions for the common point $H$ and
the scalar point $S-P$.

The required approximation is the same one used in the source.
Choose $H$ sufficiently deep at these finite places and in a
small real neighborhood of the identity, excluding the identity
itself.  Shrink the real neighborhood so that its formal
parameter is bounded away from zero.  Choose $S-P$ still deeper
at every fixed finite place and sufficiently close to the real
identity.  The reciprocal-abscissa tests at these points then
tend to zero, and the quotient of their formal parameters
tends to zero with $H$ in that smaller neighborhood.
Fixed powers of $\rho$ in the denominators do not affect this
convergence.

One can impose these conditions with the required congruences
on the integer indices: the source first chooses the fixed
modulus and depths, then applies its prime-pattern and common-point
construction.  Increasing the modulus to include the finite
support of $\delta_X$ is allowed because this support was fixed
before that application.  It introduces no new family of
varying-prime conditions and no stronger prime-pattern theorem.
The real compact-group argument permits the same shrinking of
the neighborhoods and arbitrarily large scalar index.  All
other source index and comparison choices are unchanged.

By Lemma~\ref{sn:positive-root}, these choices make
$W_\rho$ a unit at every prime in the support of $R_X$ and
give $|W_\rho-1|<1/2$ at the real place.  Therefore
\[
                       0<b=R_XW_\rho^2<9/4<17.
\]
The factor $W_\rho^2$ preserves the local squareclasses and,
at $2,3,17$, preserves the depths of Proposition~
\ref{sn:prepared-constant}.  The reduced denominator of $b$
is a square: at a prime in the support of $R_X$, the root is
a unit, while at every other prime its valuation contribution
is even.  Thus $b$ is prepared.  The same inherited rational
pole-realization theorem supplies the final pole witnesses.
This proves completeness without a correction branch.
\end{proof}

The input polynomial has degree four and can be evaluated with
two quadratic gates while preserving the source expression
$\beta=2Y+3$:
\[
                    Z=X^2,\qquad Y=108Z^2+225228.
\]
The gain in any compiled polynomial also depends on its exact
occurrence and denominator inventory; the preceding propositions
establish the arithmetic substitution independently of that count.

\subsection{Number-field scope}

All displayed algebraic identities remain valid in characteristic
zero.  Positivity of the contact and norm root also holds in a
number field with a real embedding, after applying that embedding.
The local squareclass table, however, concerns rational integer
inputs and the three specified completions of $\mathbb Q$.
It does not establish the corresponding preparation assertions
for every algebraic integer in a larger field.

Even the integrality equivalence of Lemma~\ref{sn:integral-input}
does not automatically become an equivalence for rings of
integers.  Let $x_0$ be a real fourth root of $1/108$ and
$L=\mathbb Q(x_0)$.  Then
\[
                       \beta(x_0)=450461\in\mathbb Z.
\]
But $x_0$ is not an algebraic integer, since $x_0^4=1/108$.
Thus this elementary input map
alone does not transfer a definition to $\mathcal O_L$.
No assertion that the complete rational defining polynomial
defines every $\mathcal O_L$ is made or follows from this
normalization.


\section{Asymmetric formal depth for the quartic}
\label{sec:asymmetric-depth}

The finite algebra constructions use truncated logarithms on both
families of elliptic points. The quartic can avoid the five primary
logarithms by imposing a larger formal depth. We give the local
argument and identify exactly the extra prime-selection input.
These are the asymmetric-depth statements of
\cite[Section ``Replacing five logarithms by asymmetric formal
depth'']{RepositoryV9}.

\begin{lemma}\label{lem:square-abscissa-cyclic}
Under the group hypothesis $2E(K)=\mathbb ZP$, a finite
nonidentity point of $E(K)$ with nonzero square abscissa belongs
to $\mathbb ZP$. Conversely every nonzero multiple of $P$ has
a representation
\[
             x=h^2,\qquad j^2=x^2-D,\qquad y=hj
\]
with $h,j\in K^\times$.
\end{lemma}
\begin{proof}
Since $P=2R_0$ and $P$ has infinite order, the group hypothesis
gives $E(K)=\mathbb ZR_0\oplus E[2](K)$, with four rational
two-torsion points. The abscissa map for the two-isogeny is the
homomorphism to $K^\times/K^{\times2}$ that takes $Q$ to
$[x(Q)]$, takes $(0,0)$ to $[-D]=[-1]$, and takes the identity
to $[1]$; see \cite[Chapter X, Section 4]{Silverman2009}.
It kills $2E(K)$. Its values at $R_0$, $(0,0)$, and $(d,0)$
are $[1+\sqrt2]$, $[-1]$, and $[d]$.
The first two classes are independent by their signs at the
two real embeddings. The third has odd valuation at the prime
$(d)$ above seven and is independent of both. Thus the image
has eight elements, equal to $|E(K)/2E(K)|$, and its kernel
is exactly $2E(K)$.

The duplication formula makes the abscissa of every nonzero
double a square. Choose its nonzero square root $h$ and put
$j=y/h$. Its curve equation gives $j^2=x^2-D$.
Neither $h$ nor $j$ vanishes on a nonzero cyclic multiple,
since such a point is not two-torsion.
\end{proof}

\paragraph{The logarithm truncation.}
For the curve $y^2=x^3-Dx$, put $f(Q)=1/x(Q)$ and
$z(Q)=-x(Q)/y(Q)$. Define
\[
 \begin{gathered}
 G_9(W)=\sum_{j=0}^{9}a_jW^j,\qquad
 a_0=1,\quad a_1=-\frac25,\\
 a_j=\frac{4j-5}{4j+1}a_{j-1}\quad(2\le j\le9),\qquad
 \ell_{40}(Q)=z(Q)G_9\bigl(Df(Q)^2\bigr).
 \end{gathered}
\]
These definitions specify all coefficients used by the five alternate
logarithm circuits.

To verify the formula, let $A(W)=\sum_{j\ge0}a_jW^j$, with the
same recurrence continued for every $j\ge2$. The curve equation gives
$z^2=f/(1-Df^2)$, while the invariant differential satisfies
$dx/(2y)=z\,df/(2f)$. Substitution shows that
$d(zA(Df^2))=dx/(2y)$ precisely when
\[
                 4W(1-W)A'(W)+(1+W)A(W)=1-W.
\]
The coefficient of $W^0$ gives $a_0=1$, the coefficient of $W$
gives $a_1=-2/5$, and every subsequent coefficient gives the displayed
recurrence. The unique formal logarithm with linear coefficient one
is therefore $zA(Df^2)$.

At an unramified place above $p>76176$ where $D$ is a unit and
$v_p(z(Q))=r\ge1$, one has $v_p(Df(Q)^2)=4r$.
The recurrence gives
\[
 v_p(a_j)\ge-\frac{4j+1}{p-1},
\]
and $a_j$ is $p$-integral if $4j+1<p$.
Indeed its denominator divides a product of integers $4k+1$ for
$1\le k\le j$, and the sum of the numbers of their multiples of
$p,p^2,\ldots$ is at most $(4j+1)/(p-1)$.
These inequalities show that the terms with $j\ge10$ have valuation
at least $40r$. Thus the truncation has relative precision at least
forty, as used in the comparison argument.

Write $U_0=H=[h]P$, $U_i=[u_i]P$, $V_0=H-P$, and
$V_i=U_i-Q_i$, where $Q_i=[n_i]P$. Put
\[
 L_i=z(U_i)/z(H),\qquad
 L'_i=\sigma\bigl(\ell_{40}(V_i)/\ell_{40}(V_0)\bigr).
\]
The nine order-one tests are $f(U_i)$ for $1\le i\le4$ and
$\sigma f(V_i)$ for $0\le i\le4$. The seven order-forty
tests are $f(H)$, the four $L_i-L'_i$, $L_1-\beta$, and
the direct Lucas residual. The scalar ratio retains order $31$.

\begin{lemma}\label{lem:asymmetric-soundness}
At a consumed unramified prime $p>76176$, these tests imply
$L_i\equiv n_i\pmod{p^{40}}$ at the selected place.
\end{lemma}
\begin{proof}
For a formal point with $r=v(z(Q))\ge1$, the formal logarithm
of this curve satisfies
\[
               \log_E(Q)/z(Q)\in1+p^{4r}\mathcal O.
\]
Indeed its first correction is of order four in $z$, and the
coefficient bound $v(a_j)\ge-(4j+1)/(p-1)$ ensures that later
terms have at least that valuation; the coefficients with
$4j+1<p$ are integral. The large-prime bound makes the same
inequality hold for all the remaining terms. The omitted tail
in $\ell_{40}$ similarly has relative precision forty.

The test $v(f(H))\ge40$ gives $v(z(H))\ge20$.
The full logarithm ratios are $u_i/h$ and
$(u_i-n_i)/(h-1)$. At least one of $h,h-1$ is a $p$-unit.
The corresponding ratio is integral, and the relative-unit
estimates show that its truncated ratio has the same valuation.
The cross-test transfers integrality to the other ratio.
Consequently $v(z(U_i))\ge v(z(H))\ge20$ as well.
The primary relative error is now divisible by $p^{80}$,
and the alternate error by $p^{40}$. Since the ratios are
integral, their absolute errors have the same bounds. Multiplying
the two ratio congruences by $h$ and $h-1$ and subtracting yields
the assertion. At an inert prime conjugation preserves the
valuation, so the identical argument applies.
\end{proof}

For completeness we use the following explicit weighted form of
the prime-pattern input. For distinct fixed rational coefficients
$c_j$, after the prescribed congruences and obstruction removal,
the patterns $p_j=u+c_jw^2$ with $u\in[N/2,N]$ and
$1\le w\le \tau(N)\sqrt N$ have total weight at least $c_0N^{3/2}\tau(N)$ for a fixed
$c_0>0$. Their individual weights are
$O((\log N)^k)$ for $k$ forms, and $\tau(N)$ can tend to zero
arbitrarily slowly in the diagonal choice. This is the source's
fixed-field weighted count; it is an arithmetic input here.

\begin{lemma}\label{lem:prime-order-separation}
Under this weighted input the selected primes may be chosen so
that no selected prime divides any opposite reduced group order.
\end{lemma}
\begin{proof}
Require additionally
$\tau(N)N^{1/6}/(\log N)^k\longrightarrow\infty$.
The strip $w<N^{1/3}$ has at most $O(N^{4/3})$ pairs and
weighted size $O(N^{4/3}(\log N)^k)$, negligible compared with
the retained weight of order $N\tau(N)\sqrt N$.
An admissible pattern therefore remains with
$N^{2/3}\le w^2=o(N)$.

All selected primes are asymptotic to $u$. By Hasse's bound
\cite[Theorem V.1.1]{Silverman2009}, every opposite group order
is $d_+(p_j)=p_j+O(\sqrt N)$. For large $N$ the ratio
$d_+(p_j)/p_i$ lies strictly between $1/2$ and $2$, so divisibility
would force equality. For $i\ne j$, equality contradicts
$|p_i-p_j|\ge c_*N^{2/3}$, where
$c_* =\min_{i\ne j}|c_i-c_j|>0$.
For $i=j$ it is impossible because $4\mid d_+(p_i)$ and
$p_i$ is odd. The one-form case needs only this last argument.
\end{proof}

Let $M_\pm=\operatorname{lcm}_p(d_\pm(p)/4)$ and
$\Pi=\prod_p p$. The original order and avoidance conditions give
$\gcd(M_-,M_+)=1$, $\gcd(M_-,L)=16$, and
$\gcd(M_+,L)=1$ for the previously fixed modulus $L$.
The lemma gives $\gcd(\Pi,M_+)=1$, and the primes may be
chosen beyond the factors of $L$. Thus the congruences
\[
 h\equiv0\pmod{\Pi^{19}M_-},\qquad
 h\equiv1\pmod{M_+},\qquad h\equiv h_*\pmod L
\]
are compatible when $64\mid h_*$. At each primary place the
factor $M_-$ first makes the point formal, and the additional
$p^{19}$ raises its depth from at least one to at least twenty
by the formal logarithm isometry
\cite[Theorem IV.6.4]{Silverman2009}.
The real approximation and the fixed scalar conditions are
preserved on this progression. This gives the required primary
depth while retaining the opposite depth and all completeness
conditions. The $29$- and $31$-variable constructions do not
need this weighted strengthening.

\subsection{The height inequalities for the quartic}

For the quartic use the fixed anchor
\[
 C_*=3^{26416453},\qquad S_0=C_*^2=3^{52832906}.
\]
The exact three-adic index calculation and integer comparisons of
Section~\ref{sec:three-adic-anchor-depth} give
\[
 c=|a-1|\ge2\cdot3^{13208226}>(80H)^{10},\qquad
             1\le H<2^{2093448}.
\]
The scalar order $31$ gives $h(\beta)\le Hc^{19/31}$.
Consequently
\[
 |a|\ge c/2>(40H)^{10},\qquad
 h(\beta)<c/80\le |a|/40,
 \quad\log(40H)<c/80\le |a|/40.
\]
Indeed $c^{12/31}>(80H)^{120/31}>80H$ proves the scalar
inequality. Also $c/80>(80H)^{10}/80>40H>\log(40H)$.

Put $B=\max(|n|,|a|,|t|,|b|)$ and let $F$ be the Lucas
residual. If $F\ne0$, then $|F|<5B^4$ and the order-forty
tests, the product-height identity and contact height give
\[
 B^2-6<7H(5B^4)^{19/40},\qquad
 B^{1/10}<7H5^{19/40}+6<24H.
\]
This contradicts the anchor. Thus $F=0$, and Lemma~\ref{lem:lucas-growth-elementary} gives $\log B>|a|-3$.
If $\beta\ne n$, use the weaker order-twelve consequence of
the order-forty input tests. The reduced numerator of
$\beta-n$ is at most $2e^{h(\beta)}B$, so
\[
 B^{5/12}<40H e^{19h(\beta)/12},\qquad
 \log B<\frac{12}{5}\log(40H)+\frac{19}{5}h(\beta)
       <\frac{31}{200}|a|<|a|-3.
\]
This is the second contradiction. The argument forces
$\beta\in\mathbb Z$ independently of the particular input map,
provided its scalar contacts and integer-input witnesses obey the
stated conditions. Section~\ref{sec:scalar-correction-free}
proves these conditions for the chosen fourth-power map.



\section{Quadratic circuits and the explicit quartic}
\label{sec:quartic-compression}

A quadratic circuit is a finite list of polynomial equations of total
degree at most two, with a designated input coordinate. Arithmetic
assignments and division equations are used only with their stated
nonzero domains. Each of the algebraic transformations below preserves the
rational input projection on its stated nonzero domain. The arithmetic identification of the
upstream projection with $\mathbb Z$ is a separate input.

\subsection{Direct primary ratios}

\begin{lemma}\label{lem:direct-primary-ratios}
Let $h_0,j_0,h_i,j_i$ be nonzero elements of a field, for
$1\le i\le r$. The equations
\[
 z_i j_i=-h_i\quad(0\le i\le r),\qquad
 L_i z_0=z_i\quad(1\le i\le r)
\]
may be replaced by
\[
 A_i=j_i h_0,\qquad A_iL_i=h_i j_0\quad(1\le i\le r).
\]
The two systems have identical projections onto the $h,j,L$
coordinates, and every extension in either system is unique.
They have respectively $2r+1$ and $2r$ auxiliary field coordinates,
counting the $L_i$ in both cases.
\end{lemma}
\begin{proof}
The first system uniquely determines $z_i=-h_i/j_i$ and
$L_i=h_i j_0/(j_i h_0)$. In the second, the uniquely determined
$A_i=j_i h_0$ is nonzero, and the same formula for $L_i$ follows.
All displayed equations are quadratic.
\end{proof}

Apply the lemma over $K=\mathbb Q(\sqrt2)$ to the five primary
points. In their square-abscissa charts,
$x_i=h_i^2$, $j_i^2=x_i^2-D$, $y_i=h_i j_i$, and
$z_i=-h_i/j_i$. The retained source domain makes every $h_i,j_i$
nonzero. The intermediate $z_i$ occur only in these four ratios;
the scalar block uses its own direct expression for the needed ratio.
Replacing five $z_i$ and four $L_i$ by four $A_i$ and four $L_i$
therefore saves one $K$-coordinate, hence two rational coordinates.
The lemma supplies both reconstruction maps, including the nonzero
conditions needed for their validity.

\subsection{Evaluating products in smaller blocks}

For a list $y_1,\ldots,y_n\in K$ define
\[
                         F(q)=\prod_{i=1}^{n}(1+y_iq).
\]
The $y_i$ in the detectors are squares of the field-valued tests.
For a partition $I_1\sqcup\cdots\sqcup I_s=\{1,\ldots,n\}$,
put $F_j(q)=\prod_{i\in I_j}(1+y_iq)$. Then
\begin{equation}\label{eq:block-product-identity}
                         F(q)=\prod_{j=1}^{s}F_j(q)
\end{equation}
is an identity in the polynomial ring
$\mathbb Z[y_1,\ldots,y_n,q]$.

The source has nine tests of order one and seven of order forty.
Their coefficients are shared among three detector evaluations.
For the nine tests, use three consecutive blocks of size three.
For the seven tests, use blocks of sizes three and four.
Compute the coefficients of each block once by multiplication of
its linear factors. Within each detector compute the necessary
powers of its rational weight, evaluate the block polynomials, and
multiply their values. Equation~\eqref{eq:block-product-identity}
shows that the resulting $W$ is exactly unchanged, for every
assignment of the test values and weights. No local or height
estimate is altered by this operation.

Every multiplication in $K$ is written in the fixed basis:
\[
 (a_0+a_1\sqrt2)(b_0+b_1\sqrt2)
 =(a_0b_0+2a_1b_1)+(a_0b_1+a_1b_0)\sqrt2.
\]
Each assigned component consequently has a quadratic defining
equation. The following count includes only the shared coefficient
calculations and their three evaluations; the test squares and final
product between the two order groups are unchanged.
\[
\begin{array}{c|rr|r}
\text{number of factors}&\text{full coefficients}&\text{smaller blocks}
   &\text{difference}\\\hline
9&66&54&12\\
7&49&43&6
\end{array}
\]
For example, three cubic blocks require nine field multiplications
for their shared coefficients, hence eighteen rational coordinates.
Each evaluation requires two rational powers, six component outputs,
and four component coordinates for the two products: twelve rational
coordinates. The total is $18+3\cdot12=54$.
The analogous $3+4$ partition requires sixteen shared coordinates
and nine coordinates per evaluation, giving $16+3\cdot9=43$.

These elementary circuit counts precede the final rational changes
of coordinates. Four final simplifications available in the full
coefficient presentation use coefficients absent from the smaller
blocks. The net saving in the finished system is therefore fourteen,
rather than eighteen. The exact certificates include the five
remaining simplifications, so this count does not rely on an estimate
of compiler behavior.

\subsection{Reusing fourth powers in the cubic blocks}

For three inputs $y_1,y_2,y_3$, write
\[
 (1+y_1q)(1+y_2q)(1+y_3q)=1+e_1q+e_2q^2+e_3q^3.
\]
Their coefficients may be computed from
\[
 e_1=y_1+y_2+y_3,\qquad
 e_2=y_1y_2+(y_1+y_2)y_3,\qquad
 e_3=e_2y_3-(y_1+y_2)y_3^2.
\]
If $y_3^2$ is already a coordinate, each of the last two
identities is a quadratic defining equation for one new
field-valued coordinate. The usual coefficient circuit keeps
$y_1y_2$ as a separate intermediate in addition to $e_2,e_3$.
The two displayed equations need no such intermediate. The last
identity follows by substituting the expression for $e_2$ and
canceling its second summand.

In the nine order-one tests, each of three cubic blocks can place
an alternate reciprocal-abscissa square in the third position.
The five alternate logarithm circuits already compute
\[
              w^2=(Df(V)^2)^2=D^2 f(V)^4.
\]
Consequently the required $y_3^2=\sigma(f(V)^4)$ is an affine
expression in an existing field-valued coordinate:
\[
 y_3^2=\sigma\left(\frac{w^2}{D^2}\right),\qquad
 \frac1{D^2}=\frac{193+132\sqrt2}{2401}.
\]
Division by this fixed nonzero coefficient and conjugation are
rational linear operations. Their order in this formula is
specified: the fixed inverse is multiplied before conjugation.
Each cubic block therefore saves one $K$-coordinate, or two
rational coordinates. Three such blocks save six coordinates.

\subsection{Constant pivots and the final polynomial}

Every constant-pivot elimination used in the coefficient certificate
starts from a retained equation $aZ+F=0$, $a\in\mathbb Q^\times$.
Substitute its unique value $Z=-F/a$ into the other equations.
An equation may also be changed by polynomial multiples of other
retained equations. These operations preserve the common zero
set, with unique reconstruction of $Z$. Their particular uses
are checked to leave all remaining equations quadratic.
Fixed nonzero rational rescaling and primitive integer
normalization of an equation likewise preserve its zeros.

The parity and primary-divisibility restriction is treated
separately in Theorem~\ref{pdanchor:merged-depth}: it preserves
soundness and has witnesses at every integer input, rather than
preserving every unrestricted auxiliary solution. The scalar
input and normalization have the analogous separately proved
completeness property. Combining these arithmetic steps with the
identities of this section gives the exact quartic of
Theorem~\ref{q456:polynomial}.


\section{The input circuit and fixed coefficients}
\label{sec:quartic-465}

We now combine the preceding parity and circuit changes with an input map
whose scalar contact is already a square in the three required local fields.
There is then no finite correction parameter. The normalization in this
section uses only an integer scaling factor and introduces no fixed
denominator prime.

The input map and the normalized contact are
\[
 \beta=216X^4+450459,\qquad
 r=\beta(\beta-1)(\beta-2401)(32\beta+343)(2\beta-1),
 \qquad R=\frac{102^2r}{(1+(17r)^2)^2}.
\]
Their integrality, local-square and denominator assertions were
proved in Section~\ref{sec:scalar-correction-free}.

\paragraph{Soundness and completeness.}
Use $R$ as the contact parameter throughout the existing scalar
norm tree, including its assigned denominator powers. Write its
root as $W_R$ and its pole input as $b=RW_R^2$. For rational
inputs $R>0$. The tree starts with a nonzero constant leaf, and
each norm operation is $u^2+Rv^2$; thus $W_R>0$ on its stated
rational-expression domain. In particular, $b\ne0$ requires no
additional inverse coordinate.

At every consumed scalar contact, Proposition~\ref{sn:contact-preservation}
preserves the positive odd contact valuation. The two terms of
$u^2+Rv^2$ then have opposite valuation parity, and the unchanged
minimum-valuation argument extracts the same assigned contact
orders. The source height and comparison estimates are written
in terms of $\operatorname h(\beta)$ itself, so their numerical
constants are unchanged. They imply $\beta\in\mathbb Z$;
Lemma~\ref{sn:integral-input} then implies $X\in\mathbb Z$.

For completeness, fix an integral input first. The rational number
$R$ and its numerator and denominator support are now fixed before
the prime-pattern construction is applied. The same simultaneous
approximation used in the source can make the six scalar tests
sufficiently small that $W_R$ is a unit at every prime in that
support and $|W_R-1|<1/2$ at the real place. The permitted fixed
modulus is enlarged by this finite support. No new condition on
the later varying contact primes is needed. It follows that
\[
                 0<b=RW_R^2<9/4<17.
\]
Its local squareclasses, depths at $2,3,17$, and square reduced
denominator have exactly the required properties. The inherited
prepared pole-realization theorem therefore applies with preparation
parameter zero. The common primary index can still be chosen
divisible by $64$, as required in Section~\ref{sec:primary-depth-anchor}.
This is a direct application of the original soundness and
completeness arguments to the new scalar parameter. It does not
claim that arbitrary old auxiliary solutions are preserved.

\paragraph{Quadratic circuit and variable count.}
In the 474-coordinate source, let $Y$ denote its former shifted
input-square coordinate, so all downstream occurrences of the
input parameter are $2Y+3$. Add only one coordinate $Z$ and use
\[
                        Z=X^2,\qquad
                        Y=108Z^2+225228.
\]
Thus $2Y+3=216X^4+450459$. Set all seven remaining correction
coordinates to one and delete their seven equations. The corrected
contact gate then has a constant-linear defining equation; substitute
its unique affine value and delete that coordinate and equation.
These steps save eight coordinates and add one.

After an invertible affine change in the existing raw-contact
coordinate, the normalization uses two added coordinates, for
$\delta$ and $R$, with the three quadratic equations
\[
 R\iota=1,\qquad
 \delta-1-289r^2=0,\qquad
 10404r\iota-\delta^2=0.
\]
Here $\iota$ is the already existing inverse
coordinate, now interpreted as $1/R$. The first equation is its
retained defining equation. The other two equations introduce the
normalization without a variable denominator. Preparation zero
deletes the old four preparation coordinates and three equations;
substitution of $U+b=2b$ in every remaining occurrence gives $U=b$.

The total count is therefore
\[
                       474-8+1-4+2=465.
\]
There are 435 quadratic residuals. This is the quadratic system before the final power and pole simplifications below.

\paragraph{The smaller fixed anchor.}
The parity and primary-index arguments hold for any fixed
$C_*=3^N$ and $S=C_*^2$. The exact local depth result
of Corollary~\ref{anchor3:integer-cutoff} shows that the physical
anchor, with the retained condition $c=a-1\ne0$, forces
\[
 |c|\ge2\cdot3^r\qquad\text{when}\qquad N=2r+1.
\]
The retained height constant is exactly
\[
 \mathcal H=2^{110}210^{1800}
       \left(\prod_{\substack{p\le76176\\p\ \mathrm{prime}}}p\right)^{19}
       +128.
\]
An exhaustive prime enumeration and exact integer calculation give
$7499$ primes in the product and show that $\mathcal H$ has
$2093448$ binary digits. Thus $\mathcal H<2^{2093448}$.
Using this sharper bound,
Corollary~\ref{anchor3:numerical-constant} gives the certified choice
\[
 r=13208226,\qquad N=26416453,\qquad S=3^{52832906}.
\]
It yields $|c|>(80\mathcal H)^{10}$ and hence also
$|c|^{12/31}>80\mathcal H$, as required by the inherited
height contradictions. Completeness still allows the scalar index
to be chosen with this fixed divisibility. No coordinate, residual,
or formal coefficient in $\mathbb Z[S]$ changes. Only the fixed
integer substituted for $S$ is reduced.

\section{A quartic in 456 total variables}
\label{sec:quartic456}

We give two further reductions of the quadratic system: a more economical
inverse-power calculation and a rescaling of the four pole certificates.
The rescaling is an exact equivalence on the already enforced nonzero
domain. The change of scalar order is justified by the local extraction
argument and the same global height bounds. Neither reduction introduces
an additional arithmetic hypothesis.

\subsection{Inverse powers with quadratic division equations}

\begin{lemma}\label{q456:inverse-powers}
Let $F$ be a field and $R\in F^\times$. Starting with $v_0=R^{-1}$,
the $r$ quadratic equations
\[
                         Rv_j=v_{j-1}^2\qquad(1\le j\le r)
\]
have a unique solution. It is
\[
                         v_j=R^{-(2^{j+1}-1)}.
\]
In particular, $R^{-31}$ requires four additional coordinates when
$R^{-1}$ is already retained.
\end{lemma}
\begin{proof}
Every equation has the nonzero coefficient $R$ at its new coordinate,
so it determines that coordinate uniquely. If $v_{j-1}=R^{-a}$,
then $v_j=R^{-(2a+1)}$. Starting from $a=1$ gives
$1,3,7,15,31,\ldots$, which proves the formula by induction.
\end{proof}

The scalar order may therefore be reduced from $32$ to $31$, replacing
five ordinary squaring coordinates by four division coordinates. The
two normalized ratio components become $D/R^{31}$; their weighted
squares have the exact form
\[
                         R(D/R^{31})^2=D^2/R^{61}.
\]
The odd exponent $61$ retains the valuation-parity property of the
scalar certificate. More generally, this argument works with every
positive assigned order, as the next lemma makes explicit.

\begin{lemma}\label{q456:arbitrary-scalar-order}
Let $p$ be a prime, let $R\in\mathbb Q^\times$ have positive odd
valuation $\nu=v_p(R)$, and form a finite binary expression by applying
$(a,b)\mapsto a^2+Rb^2$ to rational leaves. Suppose a leaf is
$L=D/R^m$, with $m\ge1$, and let $W$ be the root. If
\[
                              v_p(RW^2)\ge0,
\]
then
\[
                       v_p(D)\ge(m-1)\nu+1\ge m.
\]
\end{lemma}
\begin{proof}
At every internal node the two summands have valuations of opposite
parity, so they do not cancel. A leaf at depth $j$ contributes a
candidate valuation
\[
                             2^jv_p(L)+e\nu,
                         \qquad 0\le e<2^j.
\]
The root valuation is the minimum of these candidates, by induction
up the tree. Since $\nu$ is odd, the assumed inequality gives
$v_p(W)>-\nu/2$. Therefore every candidate exceeds $-\nu/2$,
and in particular
\[
 v_p(L)>-\frac{e+1/2}{2^j}\nu>-\nu.
\]
Substituting $v_p(L)=v_p(D)-m\nu$ and using the integrality of
valuations gives the result. This proof uses no parity condition on $m$.
\end{proof}

For the scalar point indices, this gives
$M_5(\beta)^{31}\mid c$, where $c=a-1\ne0$. Under the retained
contact-height input, it follows that
\[
                            h(\beta)\le H|c|^{19/31}.
\]
The existing anchor gives $|c|>(80H)^{10}$, with $H\ge1$. Hence
\[
 |c|^{12/31}>(80H)^{120/31}>80H,
 \qquad h(\beta)<|c|/80.
\]
The order-forty comparison tests, the product-height identity, and
the two Lucas height contradictions are unchanged. This is the same
scalar upper bound used there. For integer-input completeness, $R$ is
fixed before the scalar approximation conditions are chosen. The
normalized leaves can be made sufficiently small at the real places
and sufficiently deep at the finitely many fixed primes for order $31$,
just as for any fixed positive order. Thus this reduction preserves
both arithmetic implications under exactly the previous inputs.

\subsection{A smaller quadratic model for the pole certificate}

It is important to distinguish the retained scaled conic coordinate
from the dimensionless parameter sometimes denoted by $e$. We write
the retained coordinate as $E$. For a nonzero pole input $b$, the two
equations before introducing the square and product coordinates are
\begin{equation}\label{q456:old-poles}
 \begin{split}
 17b p^2&=(17b-E)(17b+3E),\\
 17b q^2&=(s^2-H_0)(s^2+3H_0),
       \qquad H_0=E(17b-b^2).
 \end{split}
\end{equation}
These are the prepared equations with preparation parameter zero.

\begin{proposition}\label{q456:pole-model}
Over any field of characteristic zero, for every $b\ne0$, the
solutions of \eqref{q456:old-poles} correspond bijectively to the
solutions in seven coordinates $u,p,v,w,h,z,F$ of the five quadratic
equations
\begin{equation}\label{q456:new-poles}
 \begin{split}
 h&=u(17-b),\\
 z&=v^2,\\
 F&=(z-h)(z+3h),\\
 17p^2-289b+(3u-34)(17u-h)&=0,\\
 17w^2-bF&=0.
 \end{split}
\end{equation}
The correspondence on the original conic coordinates is
\[
                        E=bu,\qquad s=bv,\qquad q=bw.
\]
If $b$ is an affine expression in the remaining coordinates,
all five equations have total degree at most two in the full
coordinate list.
\end{proposition}
\begin{proof}
The first rescaling gives $H_0=b^2u(17-b)=b^2h$. Dividing the
first equation of \eqref{q456:old-poles} by $b$, and the second
by $b^3$, gives
\[
 17p^2=b(17-u)(17+3u),\qquad
 17w^2=b(v^2-h)(v^2+3h).
\]
The equations defining $z,F$ express the latter relation quadratically.
For the former, put
\[
 C=17p^2-b(289+34u-3u^2),\qquad
 G=h-17u+bu.
\]
Then the exact polynomial identity
\[
 C-(3u-34)G
       =17p^2-289b+(3u-34)(17u-h)
\]
removes the cubic term, using the already retained equation $G=0$.
This proves the equivalence of the displayed systems. The maps
$u=E/b$, $v=s/b$, $w=q/b$ and their inverses are rational and
mutually inverse, while $h,z,F$ are uniquely determined. The
degree assertion follows directly from \eqref{q456:new-poles}.
\end{proof}

The preceding realization uses four base conic coordinates and only
three auxiliary coordinates. The former quadratic realization used
the same four base coordinates and five auxiliary coordinates for
$p^2,b^2,H_0,s^2,q^2$ (up to fixed nonzero rational rescalings).
It therefore saves two coordinates and two equations per pole block.
There are four blocks, so the saving is eight coordinates.

The nonzero hypothesis is already enforced in every block. In a
comparison block write $W=w_0+w_1\sqrt2$. The retained inverse of
$w_0$ makes $w_0\ne0$. Consequently
$w_0^2-2w_1^2\ne0$ over $\mathbb Q$. The equation
$2w_0T=w_0^2-2w_1^2$ then gives $T\ne0$, and the existing
equation expressing $b$ as $T^2$ divided by the weight and its
contact difference forces $b\ne0$. This is an algebraic argument,
independent of the global height hypotheses.

For the scalar block, the fourth-power input map makes its contact
$r$ positive, and the normalization
\[
                    R=\frac{102^2r}{(1+(17r)^2)^2}
\]
makes $R>0$ for every rational input. The scalar binary expression
starts with the leaf $1$, and each operation is $a^2+Rb^2$.
Its root $W$ is therefore positive. Thus its pole input $b=RW^2$
is positive, in particular nonzero. The rescaling loses no scalar
solution either.

\subsection{The resulting coefficient table}

The construction is obtained directly from the valid 474-coordinate
source. The fourth-power input, removal of the scalar correction
coordinates, zero preparation, and integral scalar normalization give
$474-8+1-4+2=465$ coordinates. The inverse-power calculation saves
one more, and the four pole rescalings save eight. The final count is
\[
                         474-8+1-4+2-1-8=456.
\]
There are 426 quadratic residual equations. Each is cleared of its
fixed rational denominators and divided by the greatest common divisor
of its integer coefficients. No parameter-dependent division is used
in this primitive normalization.

\begin{theorem}\label{q456:polynomial}
Let $R_1,\ldots,R_{426}$ be the explicitly supplied primitive quadratic
residuals, in 456 total rational coordinates including the input $X$,
and set
\[
                           Q_{456}=\sum_{i=1}^{426}R_i^2.
\]
Its zero set over $\mathbb Q$ is exactly the zero set of the displayed
quadratic system. It has degree exactly four, content one, 27,784
ordinary monomials, and coefficient height
\[
                          H(Q_{456})=64\cdot3^{105665812}.
\]
The fixed integer coefficient parameter is $S=3^{52832906}$.
Under the four remaining stated source arithmetic inputs and the
weighted prime-selection input of the quartic construction, together
with the proved Lucas estimate, the existential input projection of
this polynomial is $\mathbb Z$.
\end{theorem}
\begin{proof}
The equivalence of the single equation and the quadratic system follows
from the ordering of $\mathbb Q$. The preceding lemmas prove the new
algebraic transformations and the order-$31$ arithmetic implications.
All other point, input, local, height, and completeness arguments of
the construction remain unchanged.

The complete coefficient list is obtained by squaring the 426 residuals
and collecting terms. Every residual has degree at most two, and the
coefficient of $X^4$ in the sum is one. Thus the degree is four and
the content is one. The table has 27,791 nonzero terms over
$\mathbb Z[S]$, which collect into 27,784 ordinary coordinate monomials.
The coefficient of the square of the first translated anchor coordinate
is exactly $64S^2$. Every other coefficient is
$c_0+c_1S+c_2S^2$, where
\[
                      |c_2|\le8,
                 \qquad |c_0|+|c_1|<2^{767}<S.
\]
Its absolute value is less than $9S^2$. These bounds also show that
none of the recorded nonzero coefficients becomes zero at the fixed
integer $S$. This proves the height and monomial assertions.

The independent checker imports no producer code. It verifies every
unchanged source row in named coordinates, the four full rescaling
identities including the cubic cancellation, all inverse-power
equations and their exponents, primitive normalization, and every
coefficient of the quartic. These finite checks certify the stated
polynomial data. The arithmetic projection assertion retains precisely
the hypotheses displayed in its statement.
\end{proof}
\section{The complete coefficient table}
The computation directory supplies all quadratic residuals and every
coefficient of their collected sum of squares. The fixed coefficient
$S$ is recorded by exact exponent notation and is not a variable.
Separate producers and checkers verify the substitutions on their
nonzero domains, the coordinate inventory, every polynomial
coefficient and the exact height. These finite certificates accompany
the proofs of the algebraic transformations. The source arithmetic
inputs are not assertions certified by these computations.
\clearpage
\begin{thebibliography}{99}
\small
\bibitem{SmallVariables2026}
\emph{Small-variable polynomial constructions over the rationals},
companion paper, 7 November 2026.
\url{https://github.com/JoseBrox/diophantine-z-q_improvements/tree/main/papers/few_variables}.

\bibitem{RepositoryV9}
E.~Glazer and contributors, \emph{Diophantine definition of the integers
in the rationals}, source repository, coefficient baseline
\texttt{14e8f601fe0d9173d4a5f084aedc3e07066a7688}, 7 November 2026.
The odd-product identity, direct targets and asymmetric-depth
construction are in \texttt{docs/quartic-523-540/addition.tex}; the
repository attributes the odd-product contribution to G.~Conant,
developed with ChatGPT.
\url{https://github.com/elliotglazer/diophantine-z-q/tree/14e8f601fe0d9173d4a5f084aedc3e07066a7688}.


\bibitem{Sun2026Rational86}
Z.-W.~Sun, \emph{86 unknowns for Hilbert's tenth problem over
the rationals}, arXiv:2610.20691v2, 19 October 2026,
especially Theorem 4.1 and Sections 10.3--10.6.
\url{https://arxiv.org/abs/2610.20691v2}.


\bibitem{Voight2021}
J.~Voight, \emph{Quaternion Algebras}, Graduate Texts in
Mathematics 288, Springer, 2021, especially Proposition 14.6.7.
\url{https://jvoight.github.io/quat.html}.


\bibitem{Silverman2009}
J.~H.~Silverman, \emph{The Arithmetic of Elliptic Curves},
second edition, Graduate Texts in Mathematics 106, Springer, 2009.
\url{https://doi.org/10.1007/978-0-387-09494-6}.


\end{thebibliography}
\end{document}
